\begin{abstract}
Multi-task leaderboards rank models by aggregating their scores on many sub-datasets, and a developer can raise a model's scores on selected sub-datasets by training on related data. We ask how much such targeted investment a given model needs to rank first, and how the amount depends on the aggregation rule. We define investment in score space (sub-datasets touched $K$, total points $L_1$, largest single gain $L_\infty$) under explicit caps on every score, compute its minimum exactly under cardinal and ordinal rules, and report the cases where no plan within the caps exists as a feasibility share. On HELM MMLU (79 models, 57 subjects), with scores allowed up to 100, the median model must raise 16 subjects by 696 points in total under the mean; the same model needs 9 subjects and 47 points under plurality, 21 subjects and 112 points under the median, and 47 subjects under Borda. The order of the rules in points is stable across caps (Kendall's $\tau=0.85$), their order in subjects touched is not ($\tau=0.63$). If a model can gain at most 10 points per subject, first place is unreachable for 63\% of the models under the mean and for every model ranked 31st or lower. On the Open LLM Leaderboard (6 tasks, 4496 models) the minimum is again cheap (paired cost 0.25 of the mean, against 0.15 on HELM), but the saving of the median (0.60, against 0.16) and of plurality (0.50, against 0.07) shrinks and Borda becomes cheaper (0.37, against 0.78), so the order of the rules changes between the two leaderboards (Kendall's $\tau$ between them 0.63 in points, 0.46 in tasks touched). The cheapest rules are also the least stable under evaluation noise. Every number comes from the released code, and the assumptions on which the conclusions rest, and where they fail, are stated.
\end{abstract}
